\scanpage{22}
% archive sheet number: 21
% source lines 1-7
Soit $G$ l'image de $\unsure{G_n|\mathscr G_n}$ dans $\mathrm{Aut.\,Aff}(E)$ (on suppose maintenant que $k$ est un corps), $\Pi$ la \deleted{polyèdre}\inserted{$n$-carte} \uline{régulière} \deleted{\unsure{primitive}} \unsure{correspondante|correspondant}, On a donc une représentation \uline{régulière} de $\Pi$. \unsure{Correspondant|Correspondent}: une représentation \uline{fidèle} de $G=\mathrm{Aut}(\Pi)$ dans $\mathrm{Aut\,Aff}(E)$, \deleted{\unsure{et}} une application \uline{injective} de $\unsure{\mathrm{Drap}(\Pi)|\mathrm{Rep}(\Pi)}$ dans $\unsure{\mathrm{Drap.\,max}(E)|\mathrm{Drap}_{\max}(E)}$.

% source lines 8-10
Les applications $D_i(\Pi)\longrightarrow\deleted{\unsure{\mathrm{Esp.}}}\unsure{\mathrm{Grass.aff}_i(E)|\mathrm{Grass}_i(E)}$ sont elles aussi injectives -- en d'autres termes, si \deleted{$\unsure{\sigma,\sigma'}\in$} $g\deleted{,\unsure{g'}}\in G\subset\mathrm{Aut.\,Aff}(E)$
% source lines 11-12: cancelled wording and its interlinear addition
\deleted{(\unsure{le|G} \unsure{engendré par} $\sigma_0,\ldots,\sigma_n$) \unsure{est tel}} que \deleted{\unsure{g et g' coïncident} \inserted{\unsure{fixant le drapeau}} \unsure{en un même} $g\cdot\unsure{f_0|\mathfrak f_0},\unsure{f_1},\ldots,\unsure{f_n}$}
% source lines 13-20
$gf_i=f_i$, $g$ est-il dans le groupe engendré par $\sigma_0\,\cdots\,\sigma_{i-1}\,\widehat{\sigma_i}\,\sigma_{i+1}\,\cdots\,\sigma_n$? Ce n'est pas évident du tout -- et c'est pourquoi il n'est pas clair que des résultats sur les \deleted{\unsure{ordres}} ordres des polygones associés à la \unsure{réal.} géom \deleted{de}\inserted{\unsure{associée|envisagée}} $\Pi$, impliquent quelque chose sur les ordres des $\sigma_i\sigma_{i+1}$!
